\notechapter{N-20220710}{Vectors, Ordered Fields, and the First Axioms of Real Numbers}

\section{Vectors express geometry using a field}

A vector in the plane may be written
\[
  \vec a=(x,y)=x\mathbf i+y\mathbf j.
\]
Addition and scalar multiplication are coordinatewise:
\[
  (x,y)+(z,w)=(x+z,y+w),\qquad
  \lambda(x,y)=(\lambda x,\lambda y).
\]
The length and dot product are
\[
  |(x,y)|=\sqrt{x^2+y^2},\qquad
  (x,y)\cdot(z,w)=xz+yw.
\]
Thus the angle between two nonzero vectors satisfies
\[
  \cos\theta=\frac{\vec a\cdot\vec b}{|\vec a|\,|\vec b|}
  =\frac{xz+yw}{\sqrt{x^2+y^2}\sqrt{z^2+w^2}}.
\]
If
\[
  \vec x=(\cos\alpha,\sin\alpha),\qquad
  \vec y=(\cos\beta,\sin\beta),
\]
then
\[
  \cos\theta=\cos\alpha\cos\beta+\sin\alpha\sin\beta
  =\cos(\alpha-\beta).
\]
This gives a vector proof of the cosine subtraction formula. The coordinate expressions depend on the field of scalars: the dot product and positive length require the ordered real-number setting, whereas the vector-space axioms themselves make sense over any field. This distinction motivates separating linear algebra from the extra order and completeness assumptions of Euclidean geometry.

A plane vector \((x,y)\) packages a displacement into two real numbers.  Vector addition corresponds to composing displacements; scalar multiplication corresponds to stretching.  The algebraic laws of vectors are therefore not arbitrary axioms.  They express basic geometric operations in coordinate form.

\section{Field laws make coordinate calculations consistent}

\begin{axiom}[Field axioms]
A field is a set \(F\) with operations \(+\) and \(\cdot\) such that \((F,+)\) is an abelian group, \((F\setminus\{0\},\cdot)\) is an abelian group, and multiplication distributes over addition.
\end{axiom}


The axioms below specify the algebraic rules shared by \(\mathbb Q\), \(\mathbb R\) and other fields; they do not yet imply that square roots exist or that Cauchy sequences converge. Their role here follows the structure of MIT 18.100A \emph{Real Analysis}, lectures on the characterization of the reals and the Archimedean property.\footnote{\href{https://ocw.mit.edu/courses/18-100a-real-analysis-fall-2020/pages/lecture-notes-and-readings/}{MIT 18.100A, \emph{Real Analysis}, lecture notes}.} The real numbers form a field. Its axioms are:
\begin{enumerate}[(F1)]
  \item \(x+(y+z)=(x+y)+z\);
  \item \(x+y=y+x\);
  \item there exists \(0\) such that \(0+x=x\);
  \item for each \(x\), there is \(-x\) such that \(x+(-x)=0\);
  \item \(x(yz)=(xy)z\);
  \item \(xy=yx\);
  \item there exists \(1\neq0\) such that \(1x=x\);
  \item for \(x\neq0\), there is \(x^{-1}\) such that \(xx^{-1}=1\);
  \item \(x(y+z)=xy+xz\).
\end{enumerate}

The note emphasizes that a ``space'' in mathematics is usually a set equipped with extra structure.  A field is not merely a set; it is a set equipped with addition, multiplication, identities, inverses, and compatibility laws.

\section{Ordered fields: what algebra and order can establish}

The order axioms include:
\begin{enumerate}[(O1)]
  \item Transitivity: if \(x\le y\) and \(y\le z\), then \(x\le z\).
  \item Antisymmetry: if \(x\le y\) and \(y\le x\), then \(x=y\).
  \item Totality: for all \(x,y\), either \(x\le y\) or \(y\le x\).
  \item Compatibility with addition: \(x\le y\Rightarrow x+z\le y+z\).
  \item Compatibility with multiplication: \(0\le x\) and \(0\le y\Rightarrow0\le xy\).
\end{enumerate}

The Archimedean axiom says that for every \(x>0\) and every \(y\), there exists a positive integer \(n\) such that
\[
  nx\ge y.
\]
A typical consequence is that every interval \((0,a)\) with \(a>0\) is nonempty: choose \(n\) so large that \(n>a^{-1}\), then \(1/n<a\).

\begin{proposition}[Consequences of an ordered field]
In an ordered field, addition preserves order and multiplication by a positive element preserves order.  In particular, squares are nonnegative:
\[
  x^2\ge0.
\]
\end{proposition}


Representative exercises include:
\begin{enumerate}[(i)]
  \item \(x\ge0\) iff \(-x\le0\), and \(y>1\) implies \(0<1/y<1\).
  \item \(1>0\) and \(-1\neq1\).
  \item If \(x\le y\) and \(a\le0\), then \(ax\ge ay\).
  \item If \(a\le b\) and \(x\le y\), then \(a+x\le b+y\), with equality iff \(a=b\) and \(x=y\).
  \item If \(a<x\) for every \(a<x\), then \(x\le y\); this is a typical contradiction argument using an interval between two distinct real numbers.
  \item \(x^2\ge0\) for every real \(x\).
  \item If \(a^2<a\), then \(0<a<1\).
  \item If nonzero \(x,y\) have the same sign, then
  \[
    (x+y)^2>(x-y)^2.
  \]
\end{enumerate}

The rationals satisfy the ordered-field axioms but lack the least-upper-bound property. A concrete witness is the set \(S=\{q\in\mathbb Q:q\ge0,\ q^2<2\}\). If it had a rational supremum \(r\), the order density of \(\mathbb Q\) and continuity of \(x\mapsto x^2\) over the reals would force \(r^2=2\), which is impossible for rational \(r\). Hence an ordered field supports algebra and order without necessarily supporting the analytic limits needed for completeness.

\section{Embedding rational numbers into an Archimedean field}

The note then proves that \(\mathbb R\) contains a copy of \(\mathbb Q\).  Define for integers \(n\)
\[
  \iota(n)=\underbrace{1+\cdots+1}_{n\text{ times}}
\]
for \(n>0\), \(\iota(0)=0\), and \(\iota(n)=-\iota(-n)\) for \(n<0\).  Then define
\[
  \iota\left(\frac pq\right)=\frac{\iota(p)}{\iota(q)}.
\]
One must check that this is well-defined: if \(p/q=s/t\), then \(pt=qs\), so
\[
  \iota(p)\iota(t)=\iota(q)\iota(s),
\]
and therefore
\[
  \frac{\iota(p)}{\iota(q)}=\frac{\iota(s)}{\iota(t)}.
\]
This embedding preserves addition, multiplication, and order.

In the real numbers, the Archimedean property says no positive infinitesimal exists inside \(\mathbb R\): for every \(x\), some integer \(n\) satisfies \(n>x\).  This leads to density of rationals and helps connect arithmetic approximations with the geometry of the line.

\section{Completeness begins where rational arithmetic stops}

The final part introduces the nested interval axiom.  If
\[
  I_n=[a_n,b_n]
\]
is a descending sequence of closed intervals with real endpoints, then
\[
  I_1\supseteq I_2\supseteq I_3\supseteq\cdots
\]
implies
\[
  \bigcap_{n\ge1}I_n\neq\varnothing.
\]
The real closed intervals are compact, so nested nonempty closed intervals have nonempty intersection. When the lengths shrink to zero, the intersection is a single point. In \(\mathbb Q\) with its inherited order and metric, there is a counterexample: take nested rational-endpoint brackets \([a_n,b_n]\cap\mathbb Q\) of \(\sqrt2\) whose lengths tend to zero. Every bracket contains rational points, yet their common intersection inside \(\mathbb Q\) is empty. This distinguishes the completeness property from the ordinary order axioms, rather than merely listing it as another axiom. The ordered field of reals can be characterized by this additional completeness principle.

The supremum property, nested intervals, Cauchy completeness, and Dedekind completeness are equivalent ways of saying that the real line has no gaps.  This is the feature that turns ordered algebra into analysis.
